% Options for packages loaded elsewhere
\PassOptionsToPackage{unicode}{hyperref}
\PassOptionsToPackage{hyphens}{url}
\documentclass[
]{article}
\usepackage{xcolor}
\usepackage{amsmath,amssymb}
\setcounter{secnumdepth}{-\maxdimen} % remove section numbering
\usepackage{iftex}
\usepackage{tabularx,ragged2e}
\ifPDFTeX
  \usepackage[T1]{fontenc}
  \usepackage[utf8]{inputenc}
  \usepackage{textcomp} % provide euro and other symbols
\else % if luatex or xetex
  \usepackage{unicode-math} % this also loads fontspec
  \defaultfontfeatures{Scale=MatchLowercase}
  \defaultfontfeatures[\rmfamily]{Ligatures=TeX,Scale=1}
\fi
\usepackage{lmodern}
\usepackage{natbib}
\ifPDFTeX\else
  % xetex/luatex font selection
\fi
% Use upquote if available, for straight quotes in verbatim environments
\IfFileExists{upquote.sty}{\usepackage{upquote}}{}
\IfFileExists{microtype.sty}{% use microtype if available
  \usepackage[]{microtype}
  \UseMicrotypeSet[protrusion]{basicmath} % disable protrusion for tt fonts
}{}
\makeatletter
\@ifundefined{KOMAClassName}{% if non-KOMA class
  \IfFileExists{parskip.sty}{%
    \usepackage{parskip}
  }{% else
    \setlength{\parindent}{0pt}
    \setlength{\parskip}{6pt plus 2pt minus 1pt}}
}{% if KOMA class
  \KOMAoptions{parskip=half}}
\makeatother
\usepackage{color}
\usepackage{fancyvrb}
\newcommand{\VerbBar}{|}
\newcommand{\VERB}{\Verb[commandchars=\\\{\}]}
\DefineVerbatimEnvironment{Highlighting}{Verbatim}{commandchars=\\\{\}}
% Add ',fontsize=\small' for more characters per line
\newenvironment{Shaded}{}{}
\newcommand{\AlertTok}[1]{\textcolor[rgb]{1.00,0.00,0.00}{\textbf{#1}}}
\newcommand{\AnnotationTok}[1]{\textcolor[rgb]{0.38,0.63,0.69}{\textbf{\textit{#1}}}}
\newcommand{\AttributeTok}[1]{\textcolor[rgb]{0.49,0.56,0.16}{#1}}
\newcommand{\BaseNTok}[1]{\textcolor[rgb]{0.25,0.63,0.44}{#1}}
\newcommand{\BuiltInTok}[1]{\textcolor[rgb]{0.00,0.50,0.00}{#1}}
\newcommand{\CharTok}[1]{\textcolor[rgb]{0.25,0.44,0.63}{#1}}
\newcommand{\CommentTok}[1]{\textcolor[rgb]{0.38,0.63,0.69}{\textit{#1}}}
\newcommand{\CommentVarTok}[1]{\textcolor[rgb]{0.38,0.63,0.69}{\textbf{\textit{#1}}}}
\newcommand{\ConstantTok}[1]{\textcolor[rgb]{0.53,0.00,0.00}{#1}}
\newcommand{\ControlFlowTok}[1]{\textcolor[rgb]{0.00,0.44,0.13}{\textbf{#1}}}
\newcommand{\DataTypeTok}[1]{\textcolor[rgb]{0.56,0.13,0.00}{#1}}
\newcommand{\DecValTok}[1]{\textcolor[rgb]{0.25,0.63,0.44}{#1}}
\newcommand{\DocumentationTok}[1]{\textcolor[rgb]{0.73,0.13,0.13}{\textit{#1}}}
\newcommand{\ErrorTok}[1]{\textcolor[rgb]{1.00,0.00,0.00}{\textbf{#1}}}
\newcommand{\ExtensionTok}[1]{#1}
\newcommand{\FloatTok}[1]{\textcolor[rgb]{0.25,0.63,0.44}{#1}}
\newcommand{\FunctionTok}[1]{\textcolor[rgb]{0.02,0.16,0.49}{#1}}
\newcommand{\ImportTok}[1]{\textcolor[rgb]{0.00,0.50,0.00}{\textbf{#1}}}
\newcommand{\InformationTok}[1]{\textcolor[rgb]{0.38,0.63,0.69}{\textbf{\textit{#1}}}}
\newcommand{\KeywordTok}[1]{\textcolor[rgb]{0.00,0.44,0.13}{\textbf{#1}}}
\newcommand{\NormalTok}[1]{#1}
\newcommand{\OperatorTok}[1]{\textcolor[rgb]{0.40,0.40,0.40}{#1}}
\newcommand{\OtherTok}[1]{\textcolor[rgb]{0.00,0.44,0.13}{#1}}
\newcommand{\PreprocessorTok}[1]{\textcolor[rgb]{0.74,0.48,0.00}{#1}}
\newcommand{\RegionMarkerTok}[1]{#1}
\newcommand{\SpecialCharTok}[1]{\textcolor[rgb]{0.25,0.44,0.63}{#1}}
\newcommand{\SpecialStringTok}[1]{\textcolor[rgb]{0.73,0.40,0.53}{#1}}
\newcommand{\StringTok}[1]{\textcolor[rgb]{0.25,0.44,0.63}{#1}}
\newcommand{\VariableTok}[1]{\textcolor[rgb]{0.10,0.09,0.49}{#1}}
\newcommand{\VerbatimStringTok}[1]{\textcolor[rgb]{0.25,0.44,0.63}{#1}}
\newcommand{\WarningTok}[1]{\textcolor[rgb]{0.38,0.63,0.69}{\textbf{\textit{#1}}}}
\usepackage{longtable,booktabs,array}
\newcounter{none} % for unnumbered tables
\usepackage{calc} % for calculating minipage widths
% Correct order of tables after \paragraph or \subparagraph
\usepackage{etoolbox}
\makeatletter
\patchcmd\longtable{\par}{\if@noskipsec\mbox{}\fi\par}{}{}
\makeatother
% Allow footnotes in longtable head/foot
\IfFileExists{footnotehyper.sty}{\usepackage{footnotehyper}}{\usepackage{footnote}}
\makesavenoteenv{longtable}
\setlength{\emergencystretch}{3em} % prevent overfull lines
\providecommand{\tightlist}{%
  \setlength{\itemsep}{0pt}\setlength{\parskip}{0pt}}
\usepackage{bookmark}
\IfFileExists{xurl.sty}{\usepackage{xurl}}{} % add URL line breaks if available
\urlstyle{same}
% fallback for those not using the hyperref driver hyperxmp:
\makeatletter
\@ifundefined{xmpquote}{\newcommand{\xmpquote}[1]{#1}}{}
\makeatother
\hypersetup{
  hidelinks,
  pdfcreator={LaTeX via pandoc}}

\author{}
\date{}

\begin{document}

\section{Hindsight Agent: A Multi-Agent for Quantitative Trading System}\label{manuscript-draft-export-to-pdf-latex}

% \begin{center}\rule{0.5\linewidth}{0.5pt}\end{center}

\subsection{Abstract}\label{abstract}

Large language models are increasingly used in financial decision
pipelines, yet published agent stacks are often hard to reproduce,
rarely enforce temporal data boundaries, and seldom report net returns
after transaction costs. We present \textbf{Hindsight}, an open
LangGraph stack with structured Pydantic outputs, a
simulation end date for point-in-time data, ticker anonymization to avoid any llm bias for company names, and a fee-aware paper backtest suite.

The primary contribution is reproducible systems
engineering, auditable stage artifacts, and frozen CSV
schedules. As an illustrative case study, we evaluate Hindsights full multi-agent pipeline
on two liquid large-cap NSE listings---\textbf{RELIANCE.NS} and
\textbf{TCS.NS}---in four pre-specified regime windows (two
tickers × exogenous bear/bull labels. one deterministic
trajectory per window, no seed averaging).
\textbf{Code and artifacts:}
\url{https://github.com/RMartires/hindsight}.


\subsection{1. Introduction}\label{introduction}

Researchers are experimenting with large language models (LLMs) inside
trading workflows amid a rapid expansion of LLM-driven agent systems. Yet evaluations rarely ask how the same implementation behaves when the underlying trend differs,
rarely enforce point-in-time data access at the adapter layer, and
rarely report net returns after stated transaction costs.

\textbf{What this paper is} \textbf{Hindsight} extends the
open-source \textbf{TradingAgents} multi-agent LangGraph framework (Xiao
et al., 2024) with (i) schema-constrained stage outputs, (ii) a
simulation end date for point-in-time vendor queries, (iii) ticker anonymization, and (iv) a scriptable, fee-aware paper backtest suite tied to frozen CSV artifacts. We document how to build
and audit such a stack and provide an honest, scoped regime-split case study on two names.

\textbf{What this paper is not.} We do not claim
state-of-the-art risk-adjusted returns, regime prediction, or
statistical generalization from a broad equity universe. We do
not report multi-seed variance (each regime window is
\textbf{N=1}). Frontier empirical papers (e.g., Jeon and Lee, 2026;
Miyazaki et al., 2026) average tens of random seeds across index
universes; our evaluation scope is deliberately narrower and is labeled
as such (§4.1, Table 3).

Case-study evaluation We run the full multi-agent pipeline
on \textbf{RELIANCE.NS} and \textbf{TCS.NS}.
For each ticker we pre-specify one bear and one bull calendar windows.
Each of four regime windows starts from \textbf{\$100,000} fresh paper cash;
we compare agent net return to buy-and-hold on the same window.


\subsection{2. Related Work}\label{related-work}

\textbf{LLM agents and classic MAS.} Multi-agent frameworks orchestrate
specialist roles for trading and finance (Xiao et al., 2024; Yu et al.,
2023; Yang et al., 2023). Bădică et al.~(2025) situate this wave as a
shift from explicitly programmed BDI-style agents to emergent LLM
behaviors raising new demands for verifiable coordination and audit
trails. While \textbf{TradingAgents} (Xiao et al., 2024) provides the
foundational analyst → debate → trader graph that we extend, it does not
ship adapter-level temporal clamps, schema-persisted stage artifacts, or
a fee-aware replication harness. \textbf{AlphaAgents} (Zhao et al.,
2025) validates institution-style role separation (Fundamental,
Sentiment, Valuation) with structured debate; \textbf{QuantAgents} (Li
et al., 2025) emphasizes manager-led meetings and
simulated trading feedback both comparators for our Bull/Bear debate
loop, but neither foregrounds point-in-time vendor routing or frozen
schedule CSVs as first-class outputs.


\textbf{Methodological standards, leakage, and overfitting} are central challenges in current benchmarks. \citet{Nguyen2026Toward} survey existing LLM financial multi-agent frameworks, exposing widespread vulnerabilities to transaction cost neglect and regime-shift blindness. To establish rigorous cross-system comparisons, they propose a set of five Consolidated Minimum Standards and introduce the \textbf{Coordination Breakeven Spread} (CBS) to explicitly measure the net-of-cost alpha generated by multi-agent coordination relative to trading drag. we report a simplified fee-and-turnover breakeven in §4.5 and defer full latency CBS to future work. Bailey et al.~(2014) show that undisclosed configuration search makes strong in-sample backtests statistically
cheap to manufacture; Hindsight responds with pre-specified
windows marking different regimes, frozen artifacts, and explicit
N=1 scope rather than implied generalization.

\textbf{Point-in-time policy and anonymization.} Early LLM trading
papers often relied on prompt-level ``do not peek'' instructions. \citet{Nguyen2026Toward} treat contamination control and
net-of-cost returns as minimum standards. Hindsight implements
the former via a simulation end date at vendor routing
 and the latter via a Zerodha delivery paper
ledger. Jeon and Lee (2026) (\textbf{BlindTrade}) provide the strongest
anonymization-first defense: tickers like ``AAPL'' can trigger memorized
pretraining associations rather than fresh reasoning; they anonymize
S\&P 500 constituents daily and report Sharpe 1.40 ± 0.22 across
20 seeds. Hindsight's \textbf{\texttt{STOCK\_\#\#\#\#}} aliases (§3.5) pursue the same goal within a LangGraph replication stack

\textbf{Fine-grained tasks and coordination topology.} Miyazaki et
al.~(2026) show that fine-grained analyst tasks (e.g., explicit
ratio and indicator computations) outperform coarse ``analyze
fundamentals'' prompts under leakage-controlled TOPIX-style evaluation
with 50 trials. Our market and fundamentals analysts are
tool-bound to concrete fetches and Tier-1 indicators
(§3.2) aligning architecture with that finding without claiming
comparable universe-scale statistics. Dewansh (2026) frames
explainable stability and rationale persistence as
audit criteria for orchestrated trading agents; our Pydantic stage
outputs and wide schedule CSVs are designed to support that traceability
(outlook → stance → \texttt{final\_signal}), even though we do not
formally score policy convergence.

\textbf{Benchmarks, simulators, and live evaluation.} Qian et al.~(2026)
introduce the \textbf{Agent Market Arena (AMA)}, a lifelong multi-market
benchmark showing that \textbf{agent architecture} often matters more
than LLM backbone choice, a motivation for holding the model fixed
(\texttt{qwen/qwen3.5-flash-02-23}) and varying regime context in §4.
Papadakis et al.~(2025) (\textbf{StockSim}) define order-level
simulation with latency, slippage, and limit-order-book microstructure;
our paper ledger is candlestick-level with stated fees and \textbf{no
LOB claim}, treating StockSim-class realism as aspirational future work.

\textbf{Empirical scope vs frontier benchmarks.} Jeon and Lee (2026), Li
et al.~(2025), and Miyazaki et al.~(2026) optimize for
\textbf{statistical stress tests} across index universes, seeds, or
simulated meetings. Hindsight optimizes for \textbf{auditability and
temporal integrity} on a two-ticker, four-window case study (Table
3)---complementary, not directly comparable on seeds or universe
breadth.

\textbf{Positioning summary.} Where \textbf{TradingAgents} supplies the
graph skeleton and \textbf{BlindTrade} / Nguyen and Pham (2026)
articulate \emph{what to measure}, Hindsight implements \emph{how
to reproduce it}: simulation end dates, structured artifacts, net returns, and frozen CSV schedules that regenerate Tables 1--4 without rerunning LLM inference.


\subsection{3. Method / System}\label{method-system}

\subsubsection{3.1 Architecture}\label{architecture}

\textbf{TradingAgentsGraph}: \textbf{Analysts} → \textbf{Bull/Bear
Researchers debate} → \textbf{trader} → \textbf{risk} → \textbf{final decision}.
Core routing follows \textbf{TradingAgents} (Xiao et al., 2024); our modifications focus on
temporal data policy, schema-constrained outputs, and the backtest/evaluation pipeline.

\subsubsection{3.2 Analyst task decomposition}\label{analyst-task-decomposition}

Each analyst role is bound to specific tools and emits a
structured AnalystReport (outlook + narrative). This
supports traceable rationale from raw data to downstream debate not
coarse role labels alone.

{\small
\begin{tabular}{@{} l p{2.4cm} p{3.6cm} p{3.6cm} @{}}
\toprule
Stage & Agent & Tools & Concrete tasks \\
\midrule
Market & Market analyst &
  \texttt{get\_stock\_data}\newline
  \texttt{get\_indicators} &
Fetch OHLCV; from a fixed 13-indicator Tier-1 catalog
(50/200 SMA, 10 EMA, MACD line/signal/histogram, RSI,
Bollinger middle/upper/lower, ATR, VWMA, MFI),
select up to 8 complementary series \\
\addlinespace
Fundamentals & Fundamentals analyst &
  \texttt{get\_fundamentals}\newline
  \texttt{get\_balance\_sheet}\newline
  \texttt{get\_cashflow}\newline
  \texttt{get\_income\_statement} &
  Pull statements and issuer profile;
  analyze revenue, margins, balance-sheet trends \\
\addlinespace
News & News analyst &
  \texttt{get\_news}\newline
  \texttt{get\_global\_news} &
  Company-targeted and macro news over a rolling window \\
\addlinespace
Sentiment & Social-media analyst &
  \texttt{get\_news} &
  Sentiment and social discussion synthesis
  (company/social queries) \\
\addlinespace
Debate & Bull / Bear researchers &
  --- &
  Argue buy vs sell thesis from shared evidence \\
\addlinespace
Execution & Trader, Risk manager &
  --- &
  Propose and gate \textbf{BUY/HOLD/SELL} \\
\bottomrule
\end{tabular}
}

Structured fields (\texttt{market\_outlook},
\texttt{fundamentals\_outlook}, \texttt{bull\_implied\_stance},
\texttt{final\_signal}, etc.) are persisted in schedule CSVs for audit, supporting \textbf{explainable stability} in the sense of Dewansh (2026): a coherent trace from evidence to executable action.

\subsubsection{3.3 Tools and structured outputs}\label{tools-and-structured-outputs}

Analyst tool calls respect the simulation end date, it acts as a hard cut-off for data as it
clamps vendor queries so agents cannot access data after the as-of trade date.
each stage outputs conform to \textbf{Pydantic} schemas to get a typesafe json object.

\subsubsection{3.5 Ticker anonymization}\label{ticker-anonymization}

LLMs trained on public text encode firm-specific narratives tied to
ticker symbols. In a point-in-time backtest, that pretraining
leakage can bias decisions.

Hindsight anonymizes the issuer for all LLM-facing prompts and tool outputs:

\begin{enumerate}
\item \textbf{Deterministic ticker aliases.}
  \texttt{TickerMapper.for\_real\_ticker()} maps each real symbol to a
  stable synthetic id (\texttt{e.g. STOCK\_\#\#\#\#}).

\item \textbf{Round-trip vendor access.}
  Tool wrappers call \texttt{resolve\_ticker\_for\_vendor()} before API calls,
  then \texttt{scrub\_ticker\_text()} on responses.

\item \textbf{News scrubbing.}
  \texttt{scrub\_news\_text()} replaces multi-word proper nouns with
  \texttt{PROPER\_NOUN} placeholders (heuristic, news-only).
\end{enumerate}

\subsubsection{3.6 Backtest metrics}\label{backtest-metrics}

The paper ledger records daily agent equity after fees. Primary metrics emphasize \textbf{total
return}, \textbf{maximum drawdown (MDD)}, and \textbf{fee drag};
risk-adjusted ratios are secondary on short windows (67--111 trading
days).

\begin{tabularx}{\linewidth}{@{} l X @{}}\toprule
Metric & Definition \\
\midrule
\textbf{Total return (net)} &
  $((V_T - V_0) / V_0)$ on agent equity after fees \\
\addlinespace
\textbf{Gross return} &
  Net return plus fees as a fraction of $V_0$ \\
\addlinespace
\textbf{Fee drag} &
  Gross return minus net return \\
\addlinespace
\textbf{Alpha vs B\&H} &
  Agent net return minus buy-and-hold return on the same adjusted closes \\
\addlinespace
\textbf{Max drawdown (MDD)} &
  Peak-to-trough on agent equity;
  \textbf{B\&H MDD} on $(V_0 \times \mathrm{close}/\mathrm{close}_0)$ \\
\addlinespace
\textbf{Excess Sharpe ratio} &
  Mean daily excess return (over risk-free) / stdev, $\times \sqrt{252}$.\footnotemark \\
\addlinespace
\textbf{Sortino / Calmar} &
  Implemented; reported in Table 2 (secondary) \\
\bottomrule
\end{tabularx}
\footnotetext{Risk-free rate is an India short-rate proxy of 6.5\% p.a.,
converted to a daily rate $r_f/252$ and applied post-hoc on the frozen equity series.
Annualization uses $\sqrt{252}$.}

\subsubsection{3.7 Observability}\label{observability}
To monitor the entire agent run for one day we use Langfuse for observability into out TradingAgent.

\newpage
\subsection{4. Case Study Evaluation}\label{case-study-evaluation}

\subsubsection{4.1 Evaluation scope and inferential
limits}\label{evaluation-scope-and-inferential-limits}

\begin{tabularx}{\linewidth}{@{} l X X @{}}
\toprule
Dimension & This study & Typical frontier benchmark \\
\midrule
\textbf{Universe} &
  2 liquid NSE large-caps (RELIANCE, TCS) &
  Index constituents (e.g., S\&P 500, TOPIX 100) \\
\addlinespace
\textbf{Regime windows} &
  4 pre-specified (2 tickers $\times$ bear/bull) &
  Full sample or rolling panels \\
\addlinespace
\textbf{Stochasticity} &
  \textbf{N=1 trajectory per window} (no seed variance) &
  20--50 independent trials \\
\addlinespace
\textbf{Regime labels} &
  Exogenous from underlying B\&H ($\pm 10\%$) &
  Often endogenous or undisclosed \\
\addlinespace
\textbf{Bias control} &
  Simulation end date + Strict anonymization &
  Strict anonymization (BlindTrade) \\
\addlinespace
\textbf{Costs} &
  Stated delivery fees (net returns) &
  Often gross returns \\
\bottomrule
\end{tabularx}

We treat §4 as an \textbf{illustrative stress test} of one stack under
contrasting underlying trends---not a hypothesis test about expected
alpha. Outcomes \textbf{observed in these four windows} should not be
extrapolated to a broad universe without further study (§5).

\textbf{Pre-specification.} Calendar windows are listed in
\textbf{regime-timeranges} \textbf{before} agent backtests.
Regime labels are assigned \textbf{exogenously} from underlying
buy-and-hold return and are \textbf{independent} of internal Bull/Bear
Researcher debate roles.

\subsubsection{4.2 Regime definition and
protocol}\label{regime-definition-and-protocol}

\begin{tabular}{@{} l l @{}}
\toprule
Regime label & Rule (underlying B\&H on processed window, per ticker) \\
\midrule
\textbf{Bear} & Total return \textbf{$\leq -10\%$} \\
\textbf{Bull} & Total return \textbf{$\geq +10\%$} \\
\bottomrule
\end{tabular}

Windows \textbf{differ by ticker} (e.g., TCS bull Jun--Sep 2024;
RELIANCE bull Jan--Jun 2025). Each run uses \textbf{\$100,000} initial cash.

\medskip
{\small
\begin{tabularx}{\linewidth}{@{} l l l X @{}}
\toprule
Ticker & Regime & Processed window \\
\midrule
RELIANCE.NS & Bear & 2024-08-01 .. 2025-01-03  \\
\addlinespace
RELIANCE.NS & Bull & 2025-01-01 .. 2025-06-03 \\
\addlinespace
TCS.NS & Bear & 2024-12-02 .. 2025-04-03 \\
\addlinespace
TCS.NS & Bull & 2024-06-03 .. 2024-09-03  \\
\bottomrule
\end{tabularx}
}

\subsubsection{4.3 Positioning vs prior agent-trading systems}\label{positioning-vs-prior-agent-trading-systems}

\textbf{Table 3.} Qualitative comparison (see §4.1 for scope limits).

{\small
\begin{tabularx}{\linewidth}{@{} X X X X X @{}}
\toprule
Feature & Hindsight & BlindTrade & QuantAgents & Expert Investment Teams \\
\midrule
Primary goal &
  Reproducible PIT stack + audit trail &
  Anonymized RL/GNN policy &
  Simulated agent meetings &
  Hierarchical analyst teams \\
\addlinespace
Architecture &
  LangGraph (staged) &
  GNN + RL &
  Meeting simulation &
  Manager--analyst hierarchy \\
\addlinespace
Stochasticity reported &
  N=1 per regime window &
  20 seeds &
  Multi-year + live tracking &
  50 trials \\
\addlinespace
Universe &
  2 NSE large-caps &
  S\&P 500 (PIT) &
  Multi-asset simulation &
  TOPIX 100 \\
\addlinespace
Bias control &
  Sim.\ end date + opt.\ anonymization &
  Strict anonymization &
  Simulation feedback &
  Knowledge cutoff controls \\
\addlinespace
Costs in backtest &
  Net (Zerodha delivery) &
  Varies &
  Varies &
  Varies \\
\addlinespace
Key \S4 metrics &
  Return, MDD, fee drag &
  Sharpe (mean $\pm$ std) &
  ARR, volatility &
  Median Sharpe \\
\bottomrule
\end{tabularx}
}

\subsubsection{4.4 Results}\label{results}

\textbf{Table 1.} Agent vs buy-and-hold (frozen CSVs, \$100,000 start).
Returns and MDD are recomputed from daily \texttt{equity} and \texttt{close}.

\medskip
\noindent\textbf{(a) Returns.} Net, gross, fee drag, and alpha vs buy-and-hold.

{\footnotesize
\setlength{\tabcolsep}{4pt}
\begin{tabular}{@{} l l r r r r r r @{}}
\toprule
Ticker & Regime & Days & B\&H & Agent (net) & Gross & Fee drag & $\alpha$ vs B\&H \\
\midrule
RELIANCE & Bear & 111 & $-17.43\%$ & $-7.54\%$ & $-6.53\%$ & 1.01 pp & $+9.89$ pp \\
RELIANCE & Bull & 108 & $+15.09\%$ & $+5.56\%$ & $+6.19\%$ & 0.64 pp & $-9.53$ pp \\
TCS      & Bear &  88 & $-20.42\%$ & $-4.83\%$ & $-2.92\%$ & 1.91 pp & $+15.59$ pp \\
TCS      & Bull &  67 & $+21.86\%$ & $+19.61\%$ & $+21.80\%$ & 2.20 pp & $-2.25$ pp \\
\bottomrule
\end{tabular}
}

\medskip
\noindent\textbf{(b) Drawdown, fees, and signals.}
MDD is peak-to-trough. Fees are in dollars. B/H/S is buy/hold/sell days.

{\footnotesize
\setlength{\tabcolsep}{5pt}
\begin{tabular}{@{} l l r r r r @{}}
\toprule
Ticker & Regime & B\&H MDD & Agent MDD & Fees (\$) & B/H/S \\
\midrule
RELIANCE & Bear & $21.0\%$ & $5.11\%$ & 1,010 & 7/28/76 \\
RELIANCE & Bull & $11.0\%$ & $2.80\%$ & 635   & 9/25/74 \\
TCS      & Bear & $23.9\%$ & $6.78\%$ & 1,910 & 12/24/52 \\
TCS      & Bull &  $5.5\%$ & $1.65\%$ & 2,198 & 13/14/40 \\
\bottomrule
\end{tabular}
}

\textbf{Table 2.} Secondary risk-adjusted metrics (post-hoc on frozen
daily equity; excess Sharpe uses \textbf{6.5\% p.a.} risk-free
rate

\medskip
\noindent
Excess Sharpe is the agent or buy-and-hold equity path.
Sortino and Calmar are on the agent path only.

\medskip
{\footnotesize
\setlength{\tabcolsep}{6pt}
\begin{tabular}{@{} l l r r r r @{}}
\toprule
Ticker & Regime &
\begin{tabular}[b]{@{}r@{}}Excess Sharpe\\(agent)\end{tabular} &
\begin{tabular}[b]{@{}r@{}}Excess Sharpe\\(B\&H)\end{tabular} &
Sortino & Calmar \\
\midrule
RELIANCE & Bear & $-4.27$ & $-2.51$ & $-3.10$ & $-0.96$ \\
RELIANCE & Bull & $+0.74$ & $+1.25$ & $+3.67$ & $+1.94$ \\
TCS      & Bear & $-2.69$ & $-3.13$ & $-2.25$ & $-0.71$ \\
TCS      & Bull & $+3.86$ & $+3.27$ & $+24.22$ & $+11.88$ \\
\bottomrule
\end{tabular}
}

Short windows and discrete daily rebalancing make risk ratios
\textbf{noisy}; we interpret Table 2 alongside Table 1, not in
isolation. RELIANCE bull is positive on total return (+5.6\%) with
positive excess Sharpe (+0.74) but lags B\&H on alpha (−9.53 pp).

\textbf{Cross-ticker patterns (observed in these windows).}

\begin{itemize}
\tightlist
\item
  \textbf{Bear (2/2):} agent \textbf{loses less} than B\&H on total
  return and shows \textbf{markedly lower MDD} (e.g., RELIANCE bear:
  5.1\% vs 21.0\%). Both runs are \textbf{SELL-heavy}, consistent with
  reduced long exposure during drawdowns.
\item
  \textbf{Bull (mixed):} TCS agent \textbf{nearly tracks} B\&H (+19.6\%
  vs +21.9\%, α −2.25 pp); RELIANCE agent \textbf{lags} (+5.6\% vs
  +15.1\%, α −9.53 pp) with 74 SELL vs 9 BUY days.
\end{itemize}

\textbf{Table 4.} Selected outlook / stance shares (\% of trading days).
Full pivot in notebook.

\medskip
\noindent
Shares are percent of trading days.
Market columns are the market analyst's outlook.
Stance columns are the bull researcher's buy share and the bear researcher's sell share.

\medskip
{
\setlength{\tabcolsep}{6pt}
\begin{tabular}{@{} l l r r r r @{}}
\toprule
Ticker & Regime &
\begin{tabular}[b]{@{}r@{}}Market\\bullish\end{tabular} &
\begin{tabular}[b]{@{}r@{}}Market\\bearish\end{tabular} &
\begin{tabular}[b]{@{}r@{}}Bull stance\\buy\end{tabular} &
\begin{tabular}[b]{@{}r@{}}Bear stance\\sell\end{tabular} \\
\midrule
RELIANCE & Bear & 10.8 & 31.5 & 100.0 & 88.3 \\
RELIANCE & Bull & 24.1 & 19.4 &  94.4 & 77.8 \\
TCS      & Bear & 22.7 & 33.0 &  71.6 & 53.4 \\
TCS      & Bull & 46.3 & 14.9 &  64.2 & 44.8 \\
\bottomrule
\end{tabular}
}

Internal \textbf{bull researcher} stances can remain buy-leaning while
\textbf{final signals} stay defensive downstream trader and risk
stages drive executable actions.

\subsubsection{4.5 Transaction-cost breakeven and architectural
boundaries}\label{transaction-cost-breakeven-and-architectural-boundaries}

\textbf{Fee drag vs alpha.} Fee drag is modest (0.6--2.2 pp of initial
capital; Table 1). It does not explain RELIANCE bull
under-participation: the −9.53 pp alpha gap dominates
0.64 pp fee drag. In bear windows, defensive positioning yields
positive alpha (+9.89 pp REL, +15.59 pp TCS) despite similar
fee levels suggesting exposure management, not costs, drives
regime-split outcomes.

\textbf{Simplified breakeven spread.} Let $K$ be non-HOLD signal days
and $F$ total fees. A coarse \textbf{fee-and-turnover breakeven} is the
average daily edge required to justify rebalancing: $F / (V_0 \cdot K)$.
For RELIANCE bull ($K=83$ actionable days, $F=\$635$), this
is ≈\textbf{7.7 bps/day} far below the \textbf{88 bps/day}
implied alpha shortfall ($9.53\%/108$ days). The binding
constraint is \textbf{SELL-heavy coordination} (74 SELL days), not
delivery fees alone.

\textbf{Architectural boundary (hypothesis).} Multi-stage debate →
trader → risk can \textbf{compress drawdown} in bear windows but
\textbf{cap rally participation} in bull windows when final policy stays
defensive despite buy-leaning researcher stances (Table 4). RELIANCE
bull exemplifies this boundary; TCS bull shows the same stack can
\textbf{nearly track} B\&H when signal balance is less skewed (13 BUY,
40 SELL). We frame this as \textbf{regime × issuer interaction under one
architecture}, not a universal win/loss verdict.

We do \textbf{not} estimate latency-driven slippage (full
\textbf{Coordination Breakeven Spread} per Nguyen and Pham, 2026); that
remains future work alongside StockSim-class order-level simulation
(Papadakis et al., 2025).

\subsubsection{4.6 Figures}\label{figures}

\begin{Shaded}
\begin{Highlighting}[]
\VariableTok{MPLCONFIGDIR}\OperatorTok{=}\NormalTok{.mpl\_cache }\ExtensionTok{.venv/bin/python}\NormalTok{ scripts/generate\_paper\_figures.py}
\end{Highlighting}
\end{Shaded}

\begin{itemize}
\tightlist
\item
  \textbf{\texttt{fig\_regime\_equity\_grid}} --- 2×2 agent vs B\&H
  panels\\
\item
  \textbf{\texttt{fig\_regime\_returns\_by\_ticker}} --- total return
  bars (four runs)\\
\item
  \textbf{\texttt{fig\_regime\_signals\_by\_ticker}} --- BUY/HOLD/SELL
  counts\\
\item
  \textbf{\texttt{fig\_regime\_outlook\_bullish}} --- bullish
  analyst-outlook share heatmap
\end{itemize}

\begin{center}\rule{0.5\linewidth}{0.5pt}\end{center}

\subsection{5. Conclusion}\label{conclusion}

We presented \textbf{Hindsight 20/20}, a reproducible point-in-time
LangGraph stack with structured outputs, optional anonymization, and
fee-aware backtesting, plus a \textbf{scoped four-window case study} on
\textbf{RELIANCE.NS} and \textbf{TCS.NS}. In the \textbf{observed bear
windows}, the full pipeline shows \textbf{lower loss and drawdown} than
buy-and-hold; \textbf{bull-window outcomes are heterogeneous},
consistent with an architectural trade-off between defensive
coordination and rally participation.

\textbf{Limitations and future work.}

\begin{itemize}
\tightlist
\item
  \textbf{Inferential scope:} N=1 per regime window; no multi-seed
  variance; two tickers---not an index universe.
\item
  \textbf{Calendars:} Per-ticker windows are not synchronous; no shared
  macro calendar claim.
\item
  \textbf{Model and ablations:} Single LLM
  (\texttt{qwen/qwen3.5-flash-02-23}); anonymization on but not ablated;
  \texttt{a1}--\texttt{a3} presets not evaluated in §4.
\item
  \textbf{Metrics:} Short windows; risk ratios are secondary; no
  latency/slippage CBS; Infosys windows in \texttt{regime-timeranges}
  planned but not run.
\item
  \textbf{Future empirical work (out of scope here):} multi-seed trials,
  broader PIT universes, anonymization ablations, and full CBS with
  measured agent latency.
\end{itemize}

Regimes are \textbf{exogenous labels}, not agent predictions.
\textbf{Not claimed:} SOTA returns; regime forecasting; statistical
generalization beyond the four frozen windows.

\begin{center}\rule{0.5\linewidth}{0.5pt}\end{center}

\subsection{Contributions}\label{contributions}

\begin{enumerate}
\def\labelenumi{\arabic{enumi}.}
\tightlist
\item
  \textbf{Reproducible backtest harness and frozen artifacts} ---
  fee-aware paper ledger, schedule CSVs, replication scripts (S5 /
  M6).\\
\item
  \textbf{Point-in-time data policy and optional anonymization} ---
  simulation end date at vendor routing; deterministic ticker aliases
  (S4 / M4--M5).\\
\item
  \textbf{Structured stage outputs} --- Pydantic-constrained artifacts
  for audit and traceability (S3 / M3).\\
\item
  \textbf{Multi-agent decision graph} with tool-bound analyst
  decomposition (S1 / M1).\\
\item
  \textbf{Honest regime-conditioned case study} --- four pre-specified
  windows, documented limits, architectural-boundary interpretation (S2
  / E1).
\end{enumerate}

\textbf{Not claimed:} SOTA returns; regime prediction; multi-seed
benchmarks; ablation studies in §4.

\begin{center}\rule{0.5\linewidth}{0.5pt}\end{center}

\subsection{Acknowledgements}\label{acknowledgements}

The agent graph implementation builds on \textbf{TradingAgents} (Xiao et
al., 2024; Tauric Research, Apache License 2.0). We thank the authors
for open-sourcing the framework.

\begin{center}\rule{0.5\linewidth}{0.5pt}\end{center}

\subsection{References}\label{references}

Badica, C., Badica, A., Ganzha, M., Ivanović, M., Paprzycki, M.,
Selişteanu, D. and Wrona, Z., 2025. Contemporary Agent Technology:
LLM-Driven Advancements vs Classic Multi-Agent Systems. \emph{arXiv
preprint} arXiv:2509.02515 {[}cs.MA{]}. https://arxiv.org/abs/2509.02515

Bailey, D.H., Borwein, J.M., López de Prado, M. and Zhu, Q.J., 2014.
Pseudo-Mathematics and Financial Charlatanism: The Effects of Backtest
Overfitting on Out-of-Sample Performance. \emph{Notices of the American
Mathematical Society}, 61(5), pp.458--471.
https://doi.org/10.1090/noti1105

Dewansh, S., 2026. Policy Convergence and Explainable Stability in
Orchestrated Multi-Agent LLM Trading Frameworks. \emph{International
Journal of Advanced Engineering and Management Science}, 12(2),
pp.058--064. https://doi.org/10.22161/ijaems.122.8

Jeon, J. and Lee, H., 2026. Can Blindfolded LLMs Still Trade? An
Anonymization-First Framework for Portfolio Optimization. \emph{arXiv
preprint} arXiv:2603.17692 {[}cs.LG{]}. https://arxiv.org/abs/2603.17692

Li, X., Zeng, Y., Xing, X., Xu, J. and Xu, X., 2025. QuantAgents:
Towards Multi-agent Financial System via Simulated Trading. In
\emph{Findings of the Association for Computational Linguistics: EMNLP
2025}, pp.17438--17464, Suzhou, China.
https://doi.org/10.18653/v1/2025.findings-emnlp.945

Miyazaki, K., Kawahara, T., Roberts, S. and Zohren, S., 2026. Toward
Expert Investment Teams: A Multi-Agent LLM System with Fine-Grained
Trading Tasks. \emph{arXiv preprint} arXiv:2602.23330 {[}cs.AI{]}.
https://arxiv.org/abs/2602.23330

Nguyen, P. and Pham, T., 2026. Toward Reliable Evaluation of LLM-Based
Financial Multi-Agent Systems: Taxonomy, Coordination Primacy, and Cost
Awareness. \emph{arXiv preprint} arXiv:2603.27539 {[}cs.MA{]}.
https://arxiv.org/abs/2603.27539

Papadakis, C., Filandrianos, G., Dimitriou, A., Lymperaiou, M., Thomas,
K. and Stamou, G., 2025. StockSim: A Dual-Mode Order-Level Simulator for
Evaluating Multi-Agent LLMs in Financial Markets. \emph{arXiv preprint}
arXiv:2507.09255 {[}cs.CE{]}. https://arxiv.org/abs/2507.09255

Qian, L., Peng, X., Wang, Y., Zhang, V.J., He, H., Li, Z., Smith, H.,
Han, Y., He, Y., Li, H., Cao, Y., Yu, Y., Lopez-Lira, A., Lu, P., Nie,
J.-Y., Xiong, G., Huang, J. and Ananiadou, S., 2026. When Agents Trade:
Live Multi-Market Trading Arena for LLM Agents. In \emph{Proceedings of
the ACM Web Conference 2026}. https://doi.org/10.1145/3774904.3792821

Xiao, Y., Sun, E., Luo, D. and Wang, W., 2024. TradingAgents:
Multi-Agents LLM Financial Trading Framework. \emph{arXiv preprint}
arXiv:2412.20138 {[}q-fin.TR{]}. https://arxiv.org/abs/2412.20138

Yang, H., Liu, X.-Y. and Wang, C.D., 2023. FinGPT: Open-Source Financial
Large Language Models. \emph{arXiv preprint} arXiv:2306.06031
{[}q-fin.ST{]}. https://arxiv.org/abs/2306.06031

Yu, Y., Li, H., Chen, Z., Jiang, Y., Li, Y., Zhang, D., Liu, R., Suchow,
J.W. and Khashanah, K., 2023. FinMem: A Performance-Enhanced LLM Trading
Agent with Layered Memory and Character Design. \emph{arXiv preprint}
arXiv:2311.13743 {[}q-fin.CP{]}. https://arxiv.org/abs/2311.13743

Zhao, T., Lyu, J., Jones, S., Garber, H., Pasquali, S. and Mehta, D.,
2025. AlphaAgents: Large Language Model based Multi-Agents for Equity
Portfolio Constructions. \emph{arXiv preprint} arXiv:2508.11152
{[}q-fin.ST{]}. https://arxiv.org/abs/2508.11152

\begin{center}\rule{0.5\linewidth}{0.5pt}\end{center}

\subsection{Changelog}\label{changelog}

{\def\LTcaptype{none} % do not increment counter
\begin{longtable}[]{@{}
  >{\raggedright\arraybackslash}p{(\linewidth - 2\tabcolsep) * \real{0.4286}}
  >{\raggedright\arraybackslash}p{(\linewidth - 2\tabcolsep) * \real{0.5714}}@{}}
\toprule\noalign{}
\begin{minipage}[b]{\linewidth}\raggedright
Date
\end{minipage} & \begin{minipage}[b]{\linewidth}\raggedright
Change
\end{minipage} \\
\midrule\noalign{}
\endhead
\bottomrule\noalign{}
\endlastfoot
2026-05-29 & Expanded references and Related Work contrasts (Nguyen and
Pham, Jeon and Lee, Bailey et al., Miyazaki et al., Li et al., Zhao et
al., Papadakis et al., Qian et al., Dewansh, Bădică et al.). \\
2026-05-29 & Multi-ticker Frozen E1: TCS + completed REL bull through
2025-06-03. \\
\end{longtable}
}

\end{document}
